% ════════════════════════════════════════════════════════════════════
%  THE DISTANCE BETWEEN TWO HOMES
%  A Physicist's Education — in Experiments, Goodbyes, and Becoming
%  Saumitra S. Phatak
%
%  Latest voice-restored edition — synthesised from three earlier drafts:
%   · "The Only Constant"            (voice, refrains, the Marathi lines)
%   · "The Distance Between Two Homes" (title, the Two-Homes theme,
%                                       Survival-Mode pulled forward,
%                                       drop caps + accent, "Latest Version")
%   · "Change Is the Only Constant"  (tighter prose, chapter titles,
%                                       pull-quotes, the quantum interlude,
%                                       About the Author)
%
%  The memoir structure and book design are retained while the prose is
%  restored toward the conversational voice of the Curious Writings essays.
%
%  Self-contained. Compile twice for the contents:
%      pdflatex "The Distance Between Two Homes - Latest.tex"
%      pdflatex "The Distance Between Two Homes - Latest.tex"
% ════════════════════════════════════════════════════════════════════

\documentclass[11pt,twoside,openany]{book}

% ── Encoding ──────────────────────────────────────────────────────────
\usepackage[T1]{fontenc}
\usepackage[utf8]{inputenc}
\usepackage[english]{babel}

% ── Typography: Palatino body, Helvetica display ──────────────────────
\usepackage{mathpazo}
\usepackage[scaled=0.92]{helvet}
\usepackage{microtype}
\usepackage{setspace}\setstretch{1.10}

% ── Trim: 5.5 x 8.5 in (intimate trade/memoir size) ───────────────────
\usepackage[paperwidth=5.5in,paperheight=8.5in,
            top=0.8in,bottom=0.8in,inner=0.85in,outer=0.62in,
            headsep=12pt,footskip=24pt]{geometry}

% ── Colour ────────────────────────────────────────────────────────────
\usepackage{xcolor}
\definecolor{ink}{HTML}{20242A}
\definecolor{accent}{HTML}{9B3F2F}   % terracotta
\definecolor{soft}{HTML}{6E747C}
\color{ink}

% ── Drop caps, pull-quotes, contents, links ───────────────────────────
\usepackage{lettrine}
\usepackage{mdframed}
\usepackage{titlesec}
\usepackage{titletoc}
\usepackage{fancyhdr}
\usepackage{emptypage}
\usepackage{needspace}
\usepackage[hidelinks]{hyperref}
\hypersetup{
  pdftitle  = {The Distance Between Two Homes},
  pdfauthor = {Saumitra S. Phatak},
  pdfsubject= {A Physicist's Education — in Experiments, Goodbyes, and Becoming}
}

% ── Part styling ──────────────────────────────────────────────────────
\titleformat{\part}[display]
  {\normalfont\sffamily\centering}
  {\color{accent}\scshape\large Part \thepart}{16pt}
  {\Huge\bfseries\color{ink}}
\titlespacing*{\part}{0pt}{0pt}{20pt}

% ── Chapter styling: big terracotta numeral, title, accent rule ───────
\titleformat{\chapter}[display]
  {\normalfont\sffamily\raggedright}
  {\raggedleft\fontsize{58}{58}\selectfont\color{accent}\thechapter}{6pt}
  {\Huge\bfseries\color{ink}\raggedright}
  [\vspace{6pt}{\color{accent}\rule{\textwidth}{1pt}}]
\titlespacing*{\chapter}{0pt}{6pt}{30pt}

% Starred chapters (Interlude, Epilogue, About) — title + rule, no number
\titleformat{name=\chapter,numberless}[display]
  {\normalfont\sffamily\raggedright}
  {}{0pt}
  {\Huge\bfseries\color{ink}\raggedright}
  [\vspace{6pt}{\color{accent}\rule{\textwidth}{1pt}}]

% ── Section styling (in-chapter sub-headings) ─────────────────────────
\titleformat{\section}[hang]
  {\normalfont\sffamily\bfseries\large\color{ink}}{}{0pt}{}
\titlespacing*{\section}{0pt}{1.5em}{0.4em}

% ── Running heads ─────────────────────────────────────────────────────
\pagestyle{fancy}
\fancyhf{}
\fancyhead[LE]{\small\itshape\nouppercase{The Distance Between Two Homes}}
\fancyhead[RO]{\small\itshape\nouppercase{\leftmark}}
\fancyfoot[C]{\small\thepage}
\renewcommand{\headrulewidth}{0.3pt}
\renewcommand{\headrule}{\hbox to\headwidth{\color{soft}\leaders\hrule height \headrulewidth\hfill}}
\fancypagestyle{plain}{\fancyhf{}\fancyfoot[C]{\small\thepage}\renewcommand{\headrulewidth}{0pt}}
\renewcommand{\chaptermark}[1]{\markboth{#1}{#1}}

% ── Paragraphs ────────────────────────────────────────────────────────
\setlength{\parindent}{1.2em}
\setlength{\parskip}{0pt}
\widowpenalty=10000
\clubpenalty=10000
\emergencystretch=1em
\frenchspacing

% ── Custom commands ───────────────────────────────────────────────────
% Drop cap to open a chapter
\newcommand{\dropcap}[2]{\lettrine[lines=3,findent=2pt,nindent=4pt]{\color{accent}\textsf{\textbf{#1}}}{#2}}

% Scene break
\newcommand{\scene}{\par\vspace{0.7em}{\centering\color{accent}\rule{0.85cm}{0.6pt}\par}\vspace{0.5em}\noindent\ignorespaces}

% Custom part page: number + title + epigraph, all on one page
\newcounter{bookpartno}
\renewcommand{\thebookpartno}{\Roman{bookpartno}}
\newcommand{\bookpart}[3]{%
  \clearpage
  \refstepcounter{bookpartno}%
  \thispagestyle{empty}%
  \vspace*{0.20\textheight}%
  \begin{center}%
    {\sffamily\scshape\large\color{accent} Part \thebookpartno}\\[16pt]
    {\sffamily\bfseries\Huge\color{ink} #1}\par
    \vspace{2.2em}%
    \begin{minipage}{0.76\textwidth}\centering
      \itshape\small #2\par\vspace{0.8em}
      \normalfont\scshape\footnotesize\color{soft} #3
    \end{minipage}%
  \end{center}%
  \addcontentsline{toc}{part}{\protect\numberline{\thebookpartno}#1}%
  \vfil\newpage}

% Pull-quote for dialogue / key lines
\newmdenv[
  leftline=true,rightline=false,topline=false,bottomline=false,
  linewidth=2pt,linecolor=accent!45,
  innerleftmargin=12pt,innerrightmargin=0pt,
  innertopmargin=3pt,innerbottommargin=3pt,
  skipabove=10pt,skipbelow=10pt,backgroundcolor=white
]{pquote}

% Unnumbered chapter that still lands in the ToC and running head
\newcommand{\starchapter}[1]{\chapter*{#1}\markboth{#1}{#1}\addcontentsline{toc}{chapter}{#1}}

% ── Table-of-contents depth ───────────────────────────────────────────
\setcounter{tocdepth}{0}

\begin{document}
\frontmatter

% ── Half title ────────────────────────────────────────────────────────
\thispagestyle{empty}
\vspace*{2.6in}
\begin{center}{\sffamily\scshape\large The Distance Between Two Homes}\end{center}
\clearpage

% ── Title page ────────────────────────────────────────────────────────
\thispagestyle{empty}
\begin{center}
  \vspace*{1.0in}
  {\sffamily\bfseries\fontsize{27}{31}\selectfont The Distance\\[3pt] Between Two Homes\par}
  \vspace{0.35in}
  {\large\itshape A PhD, Two Homes,\par}
  {\large\itshape and Too Many Questions\par}
  \vfill
  {\Large Saumitra S. Phatak\par}
  \vspace{0.55in}
  {\color{accent}\rule{1.3in}{0.8pt}\par}
  \vspace{0.22in}
  {\small Mumbai \(\;\longleftrightarrow\;\) West Lafayette \(\;\longleftrightarrow\;\) Boston\par}
\end{center}
\clearpage

% ── Colophon ──────────────────────────────────────────────────────────
\thispagestyle{empty}
\vspace*{\fill}
\begin{center}\begin{minipage}{0.8\textwidth}\small
Shaped from personal essays written between 2021 and 2025. The events are true
to the author's memory; the chronology has been tightened, and some scenes have
been joined for the sake of one continuous story.\par
\vspace{1em}
Based on the \textit{Curious Writings} series.\\
Copyright \copyright\ 2026 Saumitra S. Phatak.
\end{minipage}\end{center}
\clearpage

% ── Dedication ────────────────────────────────────────────────────────
\thispagestyle{empty}
\vspace*{\fill}
\begin{center}\itshape
For Aai and Baba ---\\[0.4em]
who always knew I would come home,\\
whose hand stayed on my head\\
however far the ground beneath my feet moved.\\[1em]
And for everyone who stayed.
\end{center}
\vspace*{\fill}
\clearpage

% ── Epigraph ──────────────────────────────────────────────────────────
\thispagestyle{empty}
\vspace*{\fill}
\begin{center}
{\large\itshape Change is the only constant.}\\[1.1em]
{\small\color{soft}--- the caption on a photograph, the first week of a PhD,\\ somewhere over the Atlantic, 2021}
\end{center}
\vspace*{\fill}
\clearpage

% ── Author's Note ─────────────────────────────────────────────────────
\starchapter{Author's Note}

I started writing these stories in 2021, mostly because too many things were
happening and I needed somewhere to put them. Moving from India to the United
States for a PhD sounds like one large event. In reality it is a thousand
small events: your first terrible meal on a long flight, a bag left on a
shuttle, a photograph from home arriving exactly when you are trying to work,
an atom finally appearing on a camera, your parents looking out at American
highways for the first time.

I wrote about these moments separately. Some essays were travel updates, some
were rants, and some were basically me asking questions that I had no ability
to answer. A few were written when the event was still fresh enough to hurt.
Others came later, when nostalgia had already started editing things in my
favour. That is memory for you. It keeps a few photographs, loses the video,
and then confidently tells you what happened.

This book clubs those writings into one story. I have cleaned the chronology,
joined a few related moments, and added what I understand now. But I have
tried not to make the twenty-five-year-old version of me sound as if he knew
the ending. He did not. Most of the time he was just trying to pass a course,
fix an experiment, cook dinner, call home, and look confident while doing all
four.

The science in these pages is real, but this is not a technical book. It is
about everything that happens around the science: the confidence you lose,
the people who return it, the family that stays involved from eight thousand
miles away, and the achievements that are far too small for a CV but somehow
large enough to change you.

\vspace{1em}
\hfill\textit{Saumitra}

% ── Contents ──────────────────────────────────────────────────────────
\clearpage
\tableofcontents

% Start Arabic page numbering without book.cls forcing a blank recto page.
\clearpage
\pagenumbering{arabic}
\makeatletter
\@mainmattertrue
\makeatother

% ══════════════════════════════════════════════════════════════════════
\bookpart{Departure}{A PhD does not begin when you enter a laboratory. It begins when the familiar scenery falls away, and you find out which parts of yourself can travel.}{Mumbai \(\to\) Chicago, 2021}
% ══════════════════════════════════════════════════════════════════════

\chapter{Wuhuu, Landed in USA!}

\dropcap{T}{he first memorable thing America gave me} was a missing bag. Not
some grand view of Chicago, not a life-changing realisation at immigration.
One bag, left behind in the airport shuttle.

I had travelled from Mumbai to Delhi, Delhi to Chicago, and then another two
and a half hours through highways and cornfields to Purdue University. A
senior picked us up and took us to our apartments. I called home and gave
everyone the successful-arrival report. The journey was smooth. America had
accepted me. I had reached the place for which we had spent months planning.
Wuhuu, landed in USA!

Then I counted my luggage. One bag was still on the shuttle.

My Indian mind immediately announced that I would never see it again. I
started explaining to myself why the things inside were not that important
anyway, which is a very useful argument the brain creates only after losing
something. Unfortunately, the things were important. So that philosophy
survived for maybe three minutes.

The bag came back two days later, completely fine. I was relieved, obviously,
but also slightly embarrassed by how quickly I had declared the situation
hopeless. That became my first lesson in America:

\begin{pquote}\itshape
Everything can be fixed, and what cannot be fixed can be lived without.
\end{pquote}

At that time I was talking about luggage. Later I used the same line for
experiments, apartments, plans, and a few versions of myself. Not bad value
from one forgotten bag, no?

\scene

The bag was only the comic ending. The real departure had happened much
earlier, and strangely it happened on the short Mumbai-to-Delhi flight, not
the international one.

As soon as the plane took off, a massive junk of thoughts came into my mind:
\textit{No, I cannot do this. I cannot leave everything I have in India for
something I may or may not have in America.} This was after applying,
accepting the offer, announcing the decision, packing everything, and reaching
the plane. Very efficient timing by the brain!

I was alone, and loneliness plus a lot of free time is not an ideal
combination for me. I tried to analyse why I was feeling this way. That did
not help. Then I started listening to music, and song by song the low feeling
disappeared. Nothing outside had changed. The plane was still taking me away
from home. But my mind was no longer concentrating only on that fact.

We always complain that music distracts us while studying. We forget that the
same distraction can be used when we are concentrating too well on the wrong
things. On that small flight, I realised that our mind focuses on what we let
it focus on. If you look straight ahead, you genuinely do not know what is
happening behind your back.

By Delhi, the panic had loosened its grip.

The international flight was sixteen hours and my first one ever. The seat
was not as good as I expected and neither was the first meal. So I remembered
a rule that had helped me many times: you cannot get disappointed if you set
your expectations low. I lowered them immediately, and suddenly the flight
felt much better.

Two Indian students were sitting beside me. We began with complaints about
United Airlines, because there is no better way to start a friendship than
complaining about the same thing. Complaining alone is useless; nobody shares
the load or even reacts properly. In a group, it becomes a conversation. Soon
we were discussing our hometowns, universities, and completely different
views about life. By the end of the flight, after sleeping for only two or
three hours, we were actually praising the airline and even the food. Humans
are funny that way. We can start rewriting an experience before it has
finished.

\scene

I had expected to post the whole journey online. During an earlier Bangalore
trip, two or three hundred people practically knew what my taxi driver looked
like. Surely going to Chicago deserved more coverage, right? Surprisingly, I
did not post a single picture. The people I really wanted to tell already
knew, and they were sad that I was going. What would everybody else do with
the information? For once, living the journey felt more important than proving
that I was living it.

Chicago gave me the compulsory ``I have landed in the USA'' excitement. It
lasted until I noticed that the airport looked like an airport and Americans
were also humans. A few things were different, yes: the toilets, roads, and
especially the cars. If you look only at the cars, it feels as if nobody in
America has money problems. Look at the people and the balance returns.

After this important international observation, my new friend and I sat in
front of a McDonald's and shamelessly ate theplas packed from home while
gossiping about Indian stand-up comedy. So much for becoming global citizens.

Then came the shuttle. Chicago slowly disappeared and cornfields appeared out
of nowhere. We wondered why the vehicle felt so fast until somebody remembered
that the speed limits were now in miles per hour. I talked with the people
around me for most of the two-and-a-half-hour ride and probably prevented two
of them from sleeping. They may still be planning revenge.

Purdue did not arrive like some grand movie destination. The buildings simply
became fewer until this new place was suddenly where I lived. Looking back, I
think that was the correct entrance. A PhD may officially begin in a
laboratory, but emotionally it begins when the familiar scenery goes away and
you find out which parts of you have travelled successfully.

\chapter{Back to Survival Mode}

\dropcap{L}{ong before America,} almost every important phase of my education
started with the same feeling: the people here know more than I do, the rules
have changed, and maybe I have arrived in the wrong room.

My first school was in Mulund, where we lived with my grandparents. I used to
win almost every competition: sport, speech, academics. You can probably guess
that the competition was not world-class. No offence to them, but it is
impossible for a person like me to win that many prizes otherwise! At five we
moved to Thane, and Saraswati Secondary School was a serious step up. For the
first couple of years I went into pure study mode just to regain my
confidence.

The school was Marathi-medium until seventh standard. In eighth, mathematics
and science suddenly switched to English. Put yourself in the shoes of a
twelve-year-old boy: everything you have understood in one vocabulary now has
to move into another language while the knowledge remains intact. It was
crazy. Again, survival mode first. Learn a little every day, pass the next
test, and eventually the strange words stop feeling strange.

That pattern would repeat for the rest of my life.

Junior college took me to the other side of Thane and outside the almost
entirely Marathi crowd of school. Hindi entered my life properly. Years later,
it somehow became the language in which I wrote poetry best. After twelfth,
most of my friends chose engineering. I chose basic science, partly to follow
my father and partly because physics had made curiosity too difficult to
ignore.

That meant forty-five minutes each way on Mumbai local trains. Those journeys
made me tougher and much more culturally open. They also taught me that
personal space is not a fundamental right; sometimes it is merely a beautiful
theory. The personality I eventually carried to America was heavily seeded
during those years.

Mumbai University added a one-hour commute and even less clarity about the
future. A master's in physics does not simply land you a job, and the
university name was not going to carry me by itself. A few dedicated
professors motivated me, and through different connections I got projects at
TIFR, the first genuinely famous Indian institution associated with my name.

The commute became ninety minutes each way. A full day could run from 7:30 in
the morning to 10 at night. Still, the harder part was the crowd. These were
easily the best minds in physics I had met. They knew so much and seemed ready
to do much more. Back to survival mode.

Again: survival mode.

I tried for a year or two and never fully settled there. Then COVID happened,
terrible for the world but, strangely, a turning point in my education. Purdue
admitted me at twenty-four. Imagine still turning your life around through
education at that age, haha!

\scene

At Purdue there was no longer only the ``cream of India'' to compete with. It
was the international cream, haha. Man, that was scary during the first
semester. I had never even stayed in a hostel in India for twenty-four years,
and now I was in another country. All my labmates were American, and I was the
first international student in the group. Coursework, teaching, research,
cooking, money, stupid American weather, transport, and loneliness all became
one large assignment.

When I say survival mode, I do not mean that every day was miserable. I mean
the phase in which you stop expecting yourself to look impressive and just
learn the next necessary thing. Hold on, improve a little, ask for help, come
back tomorrow. Eventually the new place loses the power to frighten you.

That pattern had repeated my whole life: five minutes to school, then across
Thane, forty-five minutes by train, one hour, ninety minutes, and finally
eight thousand miles. Every phase first made me feel smaller and then made my
world larger. Purdue was new, sure, but maybe it was also the same old story
on a much bigger map.

\chapter{Too Much of Myself}

\dropcap{D}{uring my first two weeks} in America, I compared everything with
India: food, roads, houses, cars, manners, grocery stores, weather, family
life, and even how seriously people took traffic rules. First-time travellers
do this with great confidence after observing maybe twelve people.

But was it really about the place? Or was it about the people? I kept asking
this question as if two weeks of data were enough to generate a universal
theory. The roads were wider, the apartments quieter, the cars larger, and
the grocery stores arranged like somebody had solved a logistics problem a
little too efficiently. Still, after the first shock, a road was a road and a
store was a store. What made the place feel different was the distance between
people. In India, life leaks into other people's lives almost automatically.
Here, everybody had a boundary, a schedule, and a door that closed properly.
Very impressive, very civilised, and sometimes very lonely.

The food was bland, obviously, but Indian grocery stores had almost everything
needed to make food exactly as it was made at home. The missing ingredient was
not always a spice. It was the room full of people. At home, cooking and
eating happen around routines that overlap. Here, everybody returned at a
different time and food often became only a tummy-filler instead of a mood
setter.

This was the part nobody had warned me about. We discuss dollars, visa
appointments, flight tickets, weather, university rankings, and whether a
pressure cooker will survive the baggage limit. We do not discuss the exact
moment when dal tastes technically correct and emotionally wrong. The recipe
is there, the cumin is there, even the frozen rotis are there, but the
background noise is missing. No one is asking what time you will eat. No one
is entering the kitchen for no reason. Food is suddenly an individual
project, and boss, that is not how I had understood food for twenty-four
years.


Cars, big houses, and fancy gadgets are the things people in India talk about
with ``some day'' in their voice. Here they were almost necessities. My
professor could not believe I planned to manage without a car, even with free
buses around Lafayette. Meanwhile, family time and relationships maintained
by daily life, things we often treat as normal in India, looked like luxuries
here. Somebody had played an Uno reverse on the whole society.

I expected making friends in another country to be difficult. Turns out, you
just have to start a dialogue. At bus stops, on campus, and around the
apartment complex, I had meaningful conversations with people from six or
seven countries within two weeks. Most people are curious. Nobody wants to
start without a reason. If you can keep finding reasons, you can stay an
extrovert almost anywhere.

Our first apartment supplied enough reasons. We hosted students whose leases
had not started, and then a flea problem nearly pushed us out of our own
place. We packed again, stayed elsewhere during pest control, took help from
people, and later arranged a party for those we helped and those who helped
us. It was fun, of course, no need to mention that. Sometimes a community
forms because everybody's inconvenience becomes public at the same time.

\scene

The strangest observation was discipline. They actually follow the rules!
Who would have thought people might obey rules made to make their own lives
easier? As a boy who had criticised Indian uncles for refusing to return a
cricket ball after we entered their property without permission, the respect
for privacy here was almost shocking. The manners were so consistent that I
sometimes could not tell whether I was talking to a human or an automated
voice. Both were equally polite.

The positive side of privacy was that nobody cared. You got your space and
some more. But once society stopped judging me, I quietly took over the job.
Now it was just me judging myself all the time, and boss, that is difficult
to escape.

It sounds funny now, but it was genuinely disorienting. In India, judgement
comes as a public service. Somebody will tell you that you have become thin,
fat, dark, fair, lazy, successful, unsuccessful, arrogant, underconfident, or
marriageable, depending on what they had for breakfast. Here, silence created
freedom. Then the same silence created a mirror. If nobody is watching you,
what kind of person do you become? If nobody is forcing a routine on you, do
you actually have one? These were not philosophical questions at the time.
They were Tuesday evening questions, asked while deciding whether to cook,
study, call home, or lie on the bed and overthink professionally.

For the first time, I had too much of myself. I had always enjoyed solitude,
but only in the necessary quantity. Here there was no family room to enter, no
old friend appearing without a plan, no crowded train forcing human contact.
The habits, interests, and emotions I thought were permanent started changing.
There is no way to record your exact self on a particular day. You can access
only the latest version and live with it.

That version had to cook, shop, clean, plan, and recover without the invisible
infrastructure of home. Even enjoyment required planning. You wake up on
Saturday and now you have to plan the weekend yourself? How ridiculous no?!
A trip means schedules, messages, a car, ticket prices, and then remembering
that everyone involved is still a student. Sometimes the big plan is cancelled
before breakfast and a smaller one is negotiated by lunch.

The silence also amplified academic judgement. I was almost failing the
mid-semester exams, research was not progressing, and teaching duties never
ended. Things pile up and you notice only when they overflow. Confidence is
crucial. If you do not hold onto it in time, the only place you will be able
to ``stand and deliver'' is in a washroom!

Fortunately, I caught up and passed the courses comfortably. Research stayed
hard, but America began to feel manageable. I knew where to buy groceries,
which bus to catch, and whom to call when something went wrong. The comparison
slowly stopped being India versus America. It became the person who had left
India versus the person forming here.

That person was not always admirable. Sometimes he was brave. Sometimes he was
irritated by a dishwasher. Sometimes he was proud of surviving a week and then
defeated by laundry. Sometimes he wanted to be deeply philosophical, and five
minutes later he wanted good chaat. Maybe that is why writing helped. It
allowed all versions to sit in one room without forcing them to agree.

Places matter, sure. But maybe their biggest effect is not how different they
look. It is how much of yourself they force you to meet. Whether that is a
good thing or not, I was still figuring out.

% ══════════════════════════════════════════════════════════════════════
\bookpart{Apprenticeship}{Most days, science is not the discovery of a new fact. It is the slow elimination of every reason you cannot yet trust what the apparatus is telling you.}{West Lafayette, 2022--2023}
% ══════════════════════════════════════════════════════════════════════

\chapter{A Better Tomorrow, Booked}

\dropcap{T}{he first trip home} began weeks before I reached the airport. Just
knowing that I was going made the difficult days better. You will notice this
in adult life: Saturdays are often better than Sundays. Not because the plans
are better, but because Saturday still has the feeling of a better tomorrow.
For the last part of that semester, India was my Sunday. The ticket itself was
doing half the emotional work.

Lord COVID made sure until the end that I was not certain I could fly. Testing
windows, lab hours, airline rules, and the uncertainty of results became more
complicated than some quantum-mechanical calculations. Six hours before
departure, the result finally came back negative. Schr\"odinger's cat was
alive, at least for travel purposes.

I landed in Delhi thinking, \textit{Yessss, I am back home!} Then I noticed I
could not breathe perfectly. Pollution. Within minutes I had become the exact
NRI I would previously have mocked, the one who leaves for five months and
returns complaining about air and noise. Trust me, though, it all faded once
the familiar colony came into view. The buildings seemed to clean the air by
themselves. When you go from point A to point B, you understand how far you
went only after coming back to A.

Home received me without asking for proof of success. This may be the most
generous and most dangerous thing about parents:

\begin{pquote}\itshape
If you go out and fail at a hundred things, you still come home a superstar in
the eyes of your parents.
\end{pquote}

At Purdue, everything was measured: grades, data, teaching, competence. At
home, you could be maximally irresponsible and still get away with it. Jet lag
ate half the visit. I fell sick. I used all my negative thoughts to keep
testing negative for COVID. My friends rearranged their routines around me,
and my parents tolerated the demands of a returning child pretending to be an
adult. It is good to be selfish sometimes, no?

The first few days at home had no sophisticated structure. Wake up late, eat
something that I had described dramatically on video calls, meet a friend,
repeat the same stories with extra sound effects, sleep at a strange hour,
and then complain that the trip was going too fast. Everyone wanted one small
piece of time. I wanted to give it and also wanted to sit at home doing
nothing. This is the strange mathematics of short visits: every person is
important, every plan is reasonable, and somehow the total is still larger
than the number of days available.

I also discovered that nostalgia is not only for big landmarks. It hides in
the lift, the sound of pressure cooker whistles from other flats, the way
someone casually enters your room without scheduling a meeting, the vegetable
seller's voice, and the exact useless discussions that happen after dinner.
In America, free time had to be created. At home, it leaked everywhere. A day
could disappear without producing anything, and instead of guilt I felt
rested. Productivity culture would not approve. Good for it.

The trip was short, because every trip home is shorter from the inside.

\scene

Coming home does not restore an old saved file. Time has continued without
your supervision. Friends have new routines and stories in which you appeared
mostly through a screen. Parents are the same people and also a few months
older. Mumbai looked familiar, but I was looking at it with eyes trained by
another place. Had home changed, or had I? Obviously both, but the second
answer was more uncomfortable.

Before leaving again, I wanted an impossible arrangement: freeze one part of
the world while fast-forwarding the other. Let me grow, but do not let my
parents age. Let friends advance, but keep my place in their daily lives. Let
the PhD take years while home stays untouched. Unfortunately, time flows for
everyone. Do not listen to Einstein and his relativity; it does not help here.

Still, the first trip was a success. Leaving India had stopped feeling like
one irreversible act; I could cross the distance in both directions. Home was
no longer simply behind me and America ahead. Each place now sat at one end of
a journey, and reaching either meant leaving the other. I hoped later trips
would feel easier and nothing important would change. Nothing! That hope was
not very realistic, but it was useful while boarding.

The second departure also taught me that the airport is a dishonest place. It
looks practical: counters, gates, screens, boarding groups. But emotionally it
is a machine for compressing love into a few bad sentences. ``Take care.''
``Message after reaching.'' ``Eat properly.'' ``Don't worry.'' Everybody knows
these lines are too small for the moment, but we say them because the boarding
announcement does not wait for better poetry.

\chapter{Coherence: Broken Again}

\dropcap{B}{y the fall of 2022,} something epic happened every morning: I
woke up at 7:30 without an alarm. A digital clock sat directly across from my
bed, and irrespective of when I slept, the first sight was almost always the
same four digits. No, I had not secretly set an alarm. My biological clock had
apparently accepted the routine more completely than I had.

Chores by eight. Breakfast alternated between banana milkshake with eggs and
ginger tea with a chicken sandwich. Check the weather before choosing clothes,
because it could change fifteen or twenty degrees within a day. Stupid
American weather! Then bus, lab, science, bus back around seven, and cooking
dinner every evening, which many students found surprisingly ambitious. One
episode or a chess game, maybe an hour of study, and sleep by eleven. That was
it. Almost the same every weekday.

There was comfort in this repetition, but also some danger. When life becomes
too systematic, even a small disturbance looks like a philosophical crisis.
One failed measurement could spoil the whole evening. One good conversation
could make the same evening excellent. I had wanted independence, and now I
had received the full package: independence from family routines, from Indian
noise, from people asking where I was going, and also from the automatic
emotional support that came bundled with all those irritating questions.

The funny part was how easily this ``coherence'' broke. At ten in the morning
I might be fully ready for an ideal day of work. Boom! A photograph from home.
A message from school friends. An invitation to some campus event. Each one
immediately sounded more fun than the schedule in my head, and now all those
thoughts travelled with me through the day. Coherence gone.

On another day, when nothing in the lab worked, the same messages were
beautiful. They rescued the day. Same photograph, same friend, completely
different effect. The power of relativity is intense, and apparently Einstein
does not even have to be invited.

\scene

My first individual project tested that coherence properly. I had to build a
radio-frequency antenna for a particular transition in trapped lithium atoms,
part of a system for state-selective imaging. The description sounded
impressive. The work felt like electrical engineering handed to someone
trained in physics, and I struggled with it enough that I went to my adviser
to discuss changing groups. His answer was neither a speech nor an empty
reassurance:

\begin{pquote}\itshape
``I think you can do it. But I can help you find an easier lab, if the variety of
projects here has become too much.''
\end{pquote}

The word \textit{easier} did its work. Nobody wants to leave because somebody
offered to find an easier lab! I stayed.

Research has no reliable scale for difficulty. A problem may be hard because
you are learning, the equipment is wrong, the idea is wrong, nobody knows the
answer, or because you made a stupid mistake on Tuesday and will discover it
three weeks later. Emotionally, all five feel almost identical. I built,
measured, failed, swapped components, learned enough engineering to be
slightly dangerous, and built again. By December one prototype worked. Fancy
scientific language aside, that felt very good.

\scene

Beyond the lab, Indiana turned colour. The fall was so beautiful that I came
to believe everyone on earth should see it once: trees holding several seasons
at the same time, a bright campus bracing for the cold, skies wide enough to
make a daily walk feel briefly cinematic. Sometimes I left the building at a
random hour and walked with no purpose, and a thought would arrive with
surprising force: \textit{this was the dream.} A year earlier I had wanted to
be somewhere nobody knew me, building an identity from nothing, and now I was
living it so continuously that I had forgotten to notice. The person inside a
dream still has groceries and bad weather and unfinished work; achievement
does not glow from every angle while it is being lived. The recognition never
lasted long. It did not need to. It gave the day enough energy to keep going.

Earlier I believed that critical thinking was simply a blessing. By then I was
not so sure; sometimes the happiest person in the room seemed to be the one
able to think the least! Still, thinking was also how I turned this repetitive
life into something worth writing about. If you are underthinking, read. If
you are overthinking, write. I had heard that line somewhere and was happily
using it as permission for both.

This was also the period when I realised that writing about yourself is much
easier than writing about other people. Other people have their own versions
of events, their own privacy, and their own right to be misunderstood
accurately. But when you write about yourself, you can be unfair with full
authority! You can admit confusion, exaggerate your stupidity, defend your
choices, and then question the defence in the next paragraph. Maybe that is
why the essays became less like updates and more like arguments with myself.

\chapter{The Atom Appeared!}

\dropcap{I}{n early 2023, a winter school} in Arizona gave us a fresh start.
My senior labmate David and I spent a week learning from some of the best
researchers in the field, in a beautiful deserted state where the scientist
inside me breathed fresh air and the Indian foodie suffered daily. Fair
exchange, I suppose.

Being in that room gave me a strange confidence. It did not mean we were
already successful, but we had reached a place where important scientific
conversations were happening and had been allowed to join them. That has to
mean the road is at least roughly correct, no? It is all about getting the
approach, man. Once you have it, the rest starts falling into place.

Back at Purdue, everybody focused on the first proper result for the lab. We
were trying to trap, cool, and image one lithium atom. That sentence is short;
the apparatus took years. You do not see an atom like a tiny ball. You collect
its photons on a camera after lasers, magnetic fields, vacuum systems,
electronics, optics, and software all agree to behave at the same time. If
that sounds difficult, good, because it was.

Then the atom appeared!

We had trapped and cooled a single lithium atom from a hot cloud, for the
first time in our lab, and imaged it with the best accuracy in the world for
that system. That ought to feel good, eh? It did. The result came from my
adviser's vision and the collective perseverance of the group: years of
components, late evenings, failed alignments, rewritten code, and many boring
decisions that we now had permission to call important.

On the camera, the evidence looked small: one distribution for no atom and
another for one atom. But that separation carried the weight of the whole
apparatus. We started writing and submitted the lab's first paper by the end
of spring. Very satisfying!

\scene

After a result works, memory becomes dishonest in a useful way. Failures turn
into steps and delays become necessary lessons. The final story sounds much
cleaner than the actual research. It is worth remembering that while we were
in the middle, we had no idea whether the ending would arrive.

After the paper, each student received an individual project and we began
working in an ``Americanised'' way, everybody chasing a specialised American
dream. The research can be boring, but you describe it in fancy language to
stay motivated!

American labmates also made cultural differences measurable. My entire
undergraduate education had cost roughly \$1000 and graduate education around
\$500. One labmate, carrying loans above \$100,000, joked that after hearing
this he almost wanted to hit me. Lucky guy, he called me. Sometimes you
appreciate your own system only after watching somebody react to the numbers.

That summer I started teaching atomic physics online to master's students at
Mumbai University. I had wanted to contribute something to people on a
similar journey. Why not donate money? Well, the only considerable income of
a PhD student is knowledge; by God's grace, we are below the poverty line in
most other resources! So time and patience were what I could give.

Teaching also exposed holes in my own understanding very efficiently. People
say the best learning happens when you teach, and annoyingly they are right.
Research had shown me the satisfaction of making one atom visible. Teaching
showed me what it feels like when an idea becomes visible to another person.

\scene

Summer days began early because the sun was always competing to wake up first.
I rode my bicycle to the lab through a small town where you could still hear
nature, people waved back, and cars respected the cyclist. Who needs a car
with such a ride? Philosophically, the lab was also physically downhill from
my house. On difficult research days, life followed the geography a little
too accurately.

At twenty-seven, career questions had started sharing space with relationships,
family, and where a future should be built. Friends were marrying, returning
to India, staying in America, or making long distance work somehow. Our
generation balances two careers, two geographies, independence, family
values, ambition, and sacrifice. No pressure!

I did not have an answer. My method was only to make the best of the present:
write because I loved writing, teach on Saturdays because I loved teaching,
go to the lab because despite everything I loved research, call my people,
and stay available. Maybe that impact is what keeps us doing the silly tasks
called work and chasing the stupid pieces of paper called money, so that in
our spare time we can do a few things that make us feel alive.

The first atom was a scientific milestone. It also made one general point very
real for me: something invisible can become visible if enough scattered
efforts finally agree. We all try, and we keep trying.

% ══════════════════════════════════════════════════════════════════════
\bookpart{Entanglements}{I was not half-present in two places. I was becoming a person made by the movement between them.}{Two homes, one experiment}
% ══════════════════════════════════════════════════════════════════════

\chapter{Which One Is Home?}

\dropcap{B}{y December 2022,} my American life had improved greatly. Instead
of fixing everything from cooking to my emotional state, I could focus on
research. I had completed a course called \textit{How to Live Abroad} without
ever receiving grades, and I felt strangely accomplished. The only major
problem remaining was physics, which was at least the problem I had travelled
there to solve.

I packed gifts and cleared some mental space for new memories. On the flight
home, the World Cup final between Argentina and France was live on every
screen. Goosebumps! Different sections of the flight cheered for Messi and
Mbapp\'e, and thousands of feet above the earth, people from everywhere were
temporarily united by one football match. Sport does this without asking for
passports or explanations. Argentina won, and naturally I took it personally.

\scene

At home, ten of us in our mid-twenties decided that our demanding lives needed
a spiritual trip to Varanasi. Our parents had doubts. We also had doubts, to
be honest, but went anyway. At the Ganga ghats and Kashi Vishwanath temple, we
found practices and stories that had existed around us since childhood but
which we had never properly examined. Of course, as freshly manufactured
NRIs, we also ate everything we had missed and complained about pollution. It
takes only a few months abroad to become slightly annoying!

Varanasi did not solve our questions about careers, relationships, or the
future. It did something else. It put our confusion next to a city much older
than all of us, and suddenly every urgent decision looked a little less
urgent. Sometimes scale is enough.

\scene

The last night at home undid the composure I had kept the whole trip.

Goodbyes never get easy. I had assumed the second one would hurt less because
we had practised. Wrong. On the last evening, I could clearly imagine how my
friends' lives would develop over the next year while I participated mostly
through a screen. My family was growing and getting old at the same time. Why
do parents have to get old along with us? Everyone wonders this and nobody has
a useful answer.

I had held my tears for most of the trip, but that night I could not. We cried
together as a family and then prepared ourselves for the next year with, of
course, a smile. There was no tragedy; I was returning to an opportunity we
all valued. That is exactly why it was difficult to explain. You can be lucky
and sad at the same time. So back I went to the stupid yet opportunistic work
life, to chase my American dream until the next visit.

\scene

Living between countries changed the word \textit{home}. At first, home was
Mumbai and Indiana was simply where work happened. Then the Indiana apartment
collected routines, friends, cookware, failures, and jokes, and became a
working-home. Mumbai remained the place where I could be maximally
irresponsible, although now every visit had a departure date attached.

For a while I thought one of these had to be the real home. Why? A person can
belong to more than one place, just not in the same way. Mumbai held the
history that needed no explanation. Indiana held the future I was trying to
construct. One knew me before this ambition; the other knew me through it.

The distance could be measured in miles, flights, time zones, missed events,
and meals shown through video calls. But both places could still exist inside
one day. That is why one photograph from home could completely break my
coherence in the lab. I was not exactly half-present in two places. I was
becoming somebody shaped by travelling between them.

But let us not make this sound too beautiful. Two lives running in parallel
also mean missing things. A wedding attended through a phone. A friendship
changing shape while you are absent. Parents and grandparents ageing during a
winter you spend aligning a laser. You can fly home, but you cannot fly back
into every ordinary day that happened without you.

Nobody tells you clearly that leaving is also a slow kind of mourning. You do
not notice it every day because you are busy building the life for which you
left. Later, the missed parts become visible. The privilege is real, and so is
the price. I no longer think the task is to choose which home is real. The
task is to carry the distance without pretending it is free.

\chapter{When My Parents Entered My World}

\dropcap{T}{here are experiences} that are too large to explain efficiently.
My parents' first trip to America is one of them. So let me start the way my
mother starts every journey: Bappa Morya.

We bought the tickets almost a year in advance. My father had never travelled
outside India, and as the date came closer, every question about airports,
immigration, and the long flight increased the household anxiety. But they
trusted God, each other, and the Emirati horses pulling the palanquin. After
their small Sinbad voyage, they reached the port of Chicago.

The preparation itself was a family project. Which clothes should be packed?
How much food can legally and morally cross the ocean? How many medicines are
too many medicines? What if immigration asks something complicated? Every
small question carried a larger feeling: this was not just a vacation, it was
my parents entering the life that had been running in their imagination for
three years. Until then, America was a set of video-call backgrounds. Now it
would have weather, roads, grocery aisles, traffic lights, and a son trying
to look like he had everything under control.

I drove two hours to collect them. That detail mattered to me. On childhood
trips, I had simply followed them without knowing the route. Now they did not
know the road and I was responsible for showing it. No birthday had made me
feel as adult as driving my parents away from that airport.

We left Chicago. The skyscrapers slid behind us and the American landscape
became the one I actually knew: highway, open field, wind turbine, long
distances between buildings. My father looked back at the city skyline and
said, \textit{``Aaplya Mumbaichya imarti tichbhar vaatatil''}, ``next to
these, our Mumbai buildings would look tiny.'' Then the skyline was gone, and
there were only windmills, and he said nothing and simply looked. At the
apartment, ready-made food welcomed the ``guests.'' Then, slowly, we turned
the flat into a home.

\scene

Jet lag put them into an afternoon sleep from which they had no intention of
returning. Now it was my job to keep \textit{them} awake, organise the day,
explain appliances, and reassure them that the silence outside did not mean
something was wrong with the city.

The role reversal was funny only because it was temporary. I explained the
dishwasher with the seriousness of a technical seminar. I gave instructions
about the stove, the trash system, the laundry, and the difference between
walking distance in India and walking distance in America. My father asked
practical questions. My mother inspected whether I was actually eating
properly. Within one day, the apartment had gained extra containers, food
planning, and the invisible administrative department that mothers carry
everywhere.

They met friends, labmates, and relatives who had previously existed for them
as names or rectangles on a screen. Watching the groups meet felt like two
cinematic universes crossing over. Their presence changed the apartment.
Meals became social again. During one quiet Sunday with cousins, I forgot for
a while that I was abroad.

America seemed to have more time. In India, twenty-four hours often fail to
contain one day; in Indiana, retirement could be practised properly. The
strongest evidence was my mother going a full month without asking for Zee
Marathi. The sun rising in the west would have surprised me less.

More importantly, they saw what no phone call could show: the road to the lab,
the grocery store, the scale of the university, the friends who had helped,
and the quiet in which I had built a life. For years they had supported an
abstraction called ``PhD in America.'' Now it had rooms, faces, distances, and
a bicycle route.

\scene

A month passed at the speed of an hour. At the airport, the direction of the
separation had reversed: I was no longer leaving their world; they were
leaving mine.

Parents first teach you to imagine a world larger than the one you can see.
Showing them my world later was deeply satisfying, not because America was
some superior destination, but because the imagined future had become a place
where we could stand together.

Their visit also confirmed something simple: no matter how old we become, we
become children again around our parents. The dependence only changes shape.
As a child, my safety rested in their competence. As an adult, more and more
of my happiness was hiding inside theirs. The ground beneath us can change;
the hand above the head remains theirs.

After they left, the apartment looked exactly the same and felt completely
wrong for a few days. Then I went back to work, because what else do you do
after a meaningful goodbye? But now the two worlds were no longer separated
only by imagination. My parents had travelled the Indiana roads, sat in the
apartment, met the people, and lived inside the ordinary life for which I had
left home. The distance between our two homes was finally something we all
understood.

That is probably why the visit was not only happy. It was also clarifying.
They could now picture the loneliness accurately, not just worry about it.
They could picture the support system too, not just hope it existed. And I
could picture their courage better. For parents, sending a child away is also
a form of travel. They remain physically in the same place, but some part of
their mind keeps making the journey.

\chapter{Who Says Achievements Need to Be Big?}

\dropcap{I}{n January 2024,} one month later than planned, my online chess
rating crossed 2000. I was so thrilled that I could not sleep. It is stupid
that one random number gives you so much excitement and self-validation, but
there we were. I explained the rating system to friends who do not even care
about chess, built an unnecessary cache in their minds, and then announced the
milestone like a major scientific result. What is the point of personal goals
if you cannot make a big deal out of them with your people, right?

Around the same time, a school crush from years back confessed that she had
felt the same way. A third-year PhD student instantly became a schoolboy with
very delayed experimental confirmation, and I bragged to my boys for an
unreasonable amount of time. Who knew a PhD student could still get that
excited about schoolyard romance? The inner child does not care much about
your academic progress.

Then I got an American driver's licence. People here get one at sixteen; I got
mine at twenty-seven. Big deal, right? Actually yes. Who decides the age at
which you are allowed to be happy about something?! The licence meant I could
drive my parents from the airport. I was so relieved that I took the weekend
off, which PhD students will recognise as another achievement of similar
difficulty.

The funny thing about small achievements is that you cannot explain their
scale to outsiders. Someone hears ``driver's licence'' and thinks of a
plastic card. I saw airport pickups, grocery freedom, emergency confidence,
and one less dependency in a country designed by and for cars. Someone hears
``chess rating'' and sees a website number. I saw hundreds of games, stupid
losses, late-night calculations, and the satisfaction of finally doing what I
had promised myself. The world sees the label; you remember the struggle.

\scene

I started learning classical music that spring. Until then I had mostly sung
Bollywood songs, so classical music was completely unknown territory. A
concert at Purdue struck a chord that I genuinely did not know existed. If
music ever connects you with someone or something beyond yourself, cherish
it. Those bonds can last.

And I learned to make rotis! Trust me, it takes a full day to prepare yourself
mentally for this activity on a weekend. With serious remote coaching from my
mother, I tried. The rotis were round enough that people ate them with a
humming sound, which is the only standard that matters.

Years earlier, I loved a non-vegetarian thali from a Mumbai place called
Suwidha: chicken curry with hot rotis. Recreating something similar in
Indiana put a slice of home on the table. I slept like a baby that afternoon,
full and content.

\scene

The last new thing I did that spring was take my own religion seriously. When
you land in a foreign country you either cling to your culture harder or drift
from it entirely, and for me drifting was never an option; I had been going to
the shakha for two years, mostly for the traditional games we had played back
home, and along the way I began reading about practices I had previously
accepted without examining.

Hinduism is endlessly mocked for having too many gods, but consider that in
physics we are perfectly comfortable with multiple forms of energy, multiple
fundamental forces, a whole zoo of elementary particles. If that multiplicity
is allowed in our description of the physical world, why is it foolish to keep
several focal points for spiritual energy: Saraswati for learning, Ganesha for
beginnings, Shiva for power and transformation? This is not inconsistency. It
is flexibility. And the other charge, that we look for God in stone. Well, do
we not all keep photographs of the people we love around the house? The object
is not confused with the person; it focuses memory and relationship. An idol
is the same kind of thing: a place to gather attention each day and point it
toward your own becoming. Religious practice, in the end, should serve your
betterment, not the other way around. This was not a proof of religion by
physics. It was a way of holding my own inheritance without embarrassment.

\scene

None of these things will appear on a CV. Maybe that is precisely their value.
Graduate school trains you to measure publications, talks, awards, years, and
offers. But what about learning to cook every day, holding friendships across
time zones, asking for help, recovering from panic, noticing the fall, driving
to an airport, or building a cultural life on purpose? Those things also
happened during the PhD. Institutions just do not have a box for them.

One person's licence is routine and another's is freedom. One dinner is just
dinner; another is a reconstructed home. A random number on a chess server can
hold a year of losses. These achievements look small only when viewed from
outside.

Inside, they are evidence that life is still happening around the PhD. That
matters. A doctorate can become greedy; it can eat evenings, weekends,
confidence, and sometimes even personality. Small achievements return a piece
of the person to himself. They say: yes, you are a researcher, but you are
also someone who can sing badly, cook decently, lose at chess dramatically,
learn slowly, call friends unnecessarily, and feel proud for reasons no
committee will understand.

Maybe they do not replace ambition, and they should not. But they give meaning
to the individual doing all that ambitious work. These are the things that
help us fall forward, right?!

% ══════════════════════════════════════════════════════════════════════
\clearpage
\starchapter{Life, According to Quantum Physics}
% ══════════════════════════════════════════════════════════════════════

\dropcap{A}{t some point in the PhD}, physics stopped being only the subject I
studied and became the language in which I compared everything else. This is
probably an occupational hazard. Spend enough time thinking about tiny
particles and soon you start explaining family, attention, and decisions with
quantum mechanics! Some of these analogies are surely stretched. I am offering
them as metaphors, not scientific proofs, so bear with me and decide which
ones survive.

The danger with physics people is that we can make even a simple emotion look
like a Hamiltonian. But the temptation is real. Quantum mechanics has this
beautiful habit of refusing clean separation: observer and system, possibility
and outcome, distance and connection, silence and fluctuation. Life is not
quantum mechanical in the literal sense, obviously. Please do not build a
religion or a grant proposal from this paragraph. But sometimes a theory gives
you not only equations, but also better metaphors for confusion you already
had.

That is why I enjoyed thinking about these connections. They were not answers.
They were tools for asking better questions. If a particle can be described
only after choosing a basis, maybe people also need context before judgement.
If measurement changes a system, maybe attention is never neutral. If
uncertainty is fundamental, maybe my demand for total clarity before every
decision was scientifically rude.

\textbf{\textsf{The observer effect.}}\quad In quantum mechanics, measuring a
system can change it. In life, observation is not passive either. Every time
we interact with someone, we alter their state a little and they alter ours.
Your presence in somebody's life is participation, not just observation.

\textbf{\textsf{Collapse.}}\quad Before measurement, a quantum system may be
in a superposition of possibilities. A measurement gives you one definite
outcome. Choices feel similar: every decision turns many ``what-ifs'' into one
``what is.'' Until you decide, life is a cloud of maybes. Once you choose, you
have to live forward from that result. Scary, but also useful; imagine having
to occupy every possible life simultaneously!

\textbf{\textsf{Coherence.}}\quad A single particle holds its coherence far
longer than a system of many; scale brings noise. Our minds are no different.
The longer you can sustain one clear line of thought without scattering it,
the better the result. Focus is the superpower most people keep meaning to
develop.

\textbf{\textsf{Uncertainty.}}\quad Heisenberg's principle says you cannot
know a particle's position and momentum perfectly at the same time; chase one
and the other blurs. In life, the tighter you grip control, the less freedom
you feel. Loosen your hold on certainty, and another dimension quietly opens.

\textbf{\textsf{Resonance.}}\quad When two systems have matching frequencies,
they exchange energy efficiently. We even say two people are on the same
wavelength. But a little mismatch keeps a dialogue alive. If everybody agrees
on everything, what will they talk about? So if you want to stay social, maybe
disagree a little. Not too much, boss, just enough to continue the
conversation.

\textbf{\textsf{Entanglement.}}\quad Entangled particles have correlations
that survive distance. We have a human comparison ready: family. Across time
zones, moods and fears still travel. Obviously this is not literal quantum
entanglement, before a physicist attacks me, but emotionally the analogy is
difficult to resist.

\textbf{\textsf{The basis set.}}\quad Every quantum state can be expressed in
a chosen basis. Our opinions also sit on foundational beliefs that we may not
notice. The basis can change, though, and that is the interesting part. If you
are arguing with somebody, the conversation can go somewhere only after you
understand what their basis is. Otherwise you may both be correct in different
coordinate systems and angry for no reason.

\textbf{\textsf{Energy and evolution.}}\quad The energy of a physical system
determines how it evolves in time. Your present energy (mental, emotional)
sets where you are going next. It sounds obvious stated plainly. It stops
being obvious at two in the afternoon on a difficult Tuesday, after six hours
at a failing experiment.

\textbf{\textsf{There is no zero.}}\quad Even in a vacuum, the quantum fields
buzz; nothing is ever truly empty. There is no such thing as nothing
happening. The quiet moments are full of background activity, things being
processed, potential gathering for the next event. Stillness is not emptiness.
It is preparation.

\scene

In the end, quantum physics is definitely still about particles, so let us not
get carried away! But it also gives us a fun vocabulary for participation,
uncertainty, and connection. We make measurements, choose among possibilities,
lose coherence, find it again, and remain connected to people who are far
away. The universe does not reveal everything at once, and neither do we.

\textit{One experiment, one conversation, one moment at a time.}

% ══════════════════════════════════════════════════════════════════════
\bookpart{Becoming}{We need events that do not advance the work, in order to remember why advancing the work matters.}{Indiana \(\to\) Boston, 2024--2025}
% ══════════════════════════════════════════════════════════════════════

\chapter{The Lab Stays. We Don't.}

\dropcap{O}{n the eighth of April, 2024,} the world went dark in the middle of
the day. The solar eclipse! Trust me, knowing the science did not make it less
mesmerising. Imagine everything around you becoming dark and one beautiful
ring remaining in the sky, probably the largest ring any of us will ever see.

Our whole PhD gang watched with the astrophysics community, their telescopes
and commentary giving us a perfect view. For three minutes, some people
shouted and others tried to photograph the darkness. As an aspiring scientist,
it felt especially satisfying. The stars literally aligned and reminded us
that the science we study is describing a real universe. These are the moments
that make the PhD journey feel worthwhile.

\scene

That summer, the senior-most student became the first person to graduate from
our lab. He had been one of the first American students I spoke with in 2021,
and his welcome, along with our postdoc's, was part of why I joined the group.
We looked back over three years of research and enjoyed the luxury of
nostalgia while still standing inside the lab. Then he said:

\begin{pquote}\itshape
``It's so weird that after five years, you move on, but the lab stays right where
it is.''
\end{pquote}

Such an interesting thought! The optical table, cables, vacuum chamber, and
unfinished tasks wait in the same places every morning. The people pass
through. A project that once occupied your entire life becomes a paper, a
folder, or an apparatus inherited by somebody who did not see it begin.
Memories are photographs, not videos. We keep a few frames and call them five
years.

\scene

Around the same time, I attended two American weddings for labmates. The
ceremonies looked familiar because movies and television had trained us well,
but the two weddings were completely different. One was carefully arranged,
outdoor and child-free, designed as a picture-perfect memory. The other was
religious and relaxed, with families and children doing their own things.
Both had websites and gift registries. You clicked one selected gift and it
arrived at the couple's house with your message. Such an efficient algorithm,
indeed reflecting the American way!

India also won the Cricket World Cup that summer. It happened on a Saturday
morning here, and the whole weekend felt successful before it had properly
begun. Most Indians slept better after that victory. Our current heroes had
finally lifted the trophy after enough heartbreaks, and the inner child was
very satisfied.

An eclipse, two weddings, a cricket victory, and a graduation do not form a
clean scientific dataset. But all of them interrupted work and reminded me
why life cannot be measured only by how much the project advanced. The
universe worked without my troubleshooting. Friends built lives beyond the
lab. A childhood dream in cricket came true. One student left while the
apparatus remained.

One day I would also leave and somebody else would inherit the urgent problem
on my table. When is the correct time to look back and enjoy the nostalgia?
Before leaving, during the goodbye, or years later? I do not know. Everyone
finds that answer separately. The lab stays. We do not. That does not make the
years less meaningful; maybe it is why we should notice them while they are
still happening.

\chapter{Boston! Finally, a Job}

\dropcap{B}{oston arrived} with fall colours and, for the first time in my
life at twenty-eight, a job! Technically it was an internship, but after
nearly a decade as a student I was going to call it a job. I moved from
academia toward industry, from small-town Indiana to a proper city, and from
studying quantum systems toward helping build machines that use them.

For the first week, while my work authorisation was processing and I was
``legally unpaid,'' I explored the river by every possible mode: walking,
running, and biking. The sunsets reflected in the water with a skyline behind
them and gave a strong Marine Drive feeling to the Mumbaikar inside me.

This mattered more than I expected. West Lafayette had been kind to me, but it
was not a city in the way Mumbai had trained me to understand a city. Boston
had people on foot, public transport, water, noise, restaurants, tourists,
students, old buildings, and enough strangers to make anonymity feel normal
again. After years of needing a reason to go anywhere, I could simply step
outside and let the place supply the reason. That is one of the great
services a city provides.

Public transport was another pleasant surprise. Tap-and-pay subway and buses
connected even places twenty-five miles away. After living in parts of America
where you are finished without a car, this felt refreshing. People say Boston
has a European vibe; half of that is surely the transport, old houses, and
streetlamps. A proper city makes everything expensive, including people's
time, but it also gives you other lives to look at while living your own.

Setting up a new apartment is painful at any age. You plan a hundred things,
order a thousand, and by day five become semi-functional. Luckily, I had
friends and family nearby. My cousin lived four blocks away, which produced
many home-cooked meals. There were hikes, a Marathi play, chess with
colleagues, and busy weekends. My adviser once joked that I have more
connections in America than he does. Boston provided supporting data.

The Marathi play was a small but important detail. Sitting in an American city
and hearing Marathi on stage does something strange to your internal map. You
realise culture is not only tied to geography; it travels through people,
WhatsApp groups, community halls, potluck boxes, and someone loudly asking
whether everybody has eaten. Suddenly Boston was not just Boston. It had a
small Maharashtra folded inside it.

\scene

At three in the morning on a Sunday, my phone rang. I kept the volume high
because living alone had trained me to treat availability as a safety system.
It was my roommate, calling from elsewhere in the duplex, whispering:

\begin{pquote}\itshape
``There's someone in our house. Maybe a homeless person. I think we should check.''
\end{pquote}

Man! I hope you never experience the chills I got in that moment. I stepped
into the hallway and asked how he was so sure. He had seen a large man looking
into our rooms, and none of the bedroom doors were locked.

We debated calling the police. My roommate, braver or stupider than me,
insisted on checking downstairs. We took a rod and flashlight and walked down
slowly. The back door was half open. Nobody was there. Only a few drops of
blood remained on the kitchen floor.

We locked everything, messaged the others, and told the landlord. By the next
day, the likely explanation was that another roommate, returning in an altered
state, had left the door open. We changed every lock. I installed a camera and
put rods in the window tracks.

Still, fear took longer to leave than the intruder, if there had been one at
all. For a student abroad, this is a special nightmare: the people you would
normally call first are thousands of miles away. Later you say, ``Whatever
doesn't kill you makes you stronger.'' In that moment, trust me, the sentence
is completely useless. It remains one hell of an experience to not remember.

\scene

The work itself was extraordinary. I worked more than ten hours on many days,
but it did not always feel exhausting because it aligned so closely with my
PhD research: quantum computing, building machines based on principles I had
spent years studying. Once again I had entered a room full of people with
frightening expertise. School, university, TIFR, Purdue, and now industry:
survival mode knew the route by then.

The quantum industry had a funny mixture of confidence and uncertainty.
Scientists and investors talked about a revolution, while nobody knew which
architecture, company, or timeline would survive. Maybe the dream would burst
like a bubble. Maybe it would transform the world like AI after COVID. Nobody
really knew, including the people paid to build it.

The PhD had not taught me to predict the future. It had taught me to work
while the answer was unavailable. So without demanding certainty, let us ride
the tide, shall we?

\chapter{Change Is the Only Constant}

\dropcap{A}{t the edge of completing the PhD,} my thoughts arrived with more
diversity than ever: nostalgia, fulfilment, excitement, gratitude, and some
confusion, obviously. But the strongest one was familiar: \textit{change is
the only constant.} I had used those words as an Instagram caption during my
first week in the United States. Back then they sounded like a nice sentence.
After the PhD, they felt experimentally verified.

Every stage had first disturbed my confidence. Mulund to Thane. Marathi to
English in the middle of school. Junior college, crowded trains, bachelor's,
master's, TIFR, Purdue, and industry. At every transition, the talent around
me appeared to increase and everybody seemed to own some vocabulary or
certainty that I lacked. Back to survival mode each time.

This is not an argument that everybody should suffer permanently. Survival
mode is the phase between two comfortable versions of yourself. Your old
methods no longer work perfectly and the new ones are not yet natural. So you
hold your competitive spirit through the high tide, learn one thing, ask one
question, and return the next day. Small improvements compound quietly. Years
later, people may call the result talent and completely miss the panic that
helped build it.

\scene

Nearly thirty years of education is longer than anybody would recommend. Was
it worth it? You decide. If education is only a transaction for employment,
my route was spectacularly inefficient: long commutes, financial limitations,
exams, anxiety, difficult research, and a doctorate completed around thirty.
There were certainly shorter ways to start earning.

But education changed much more than my employment. Marathi-medium school
became a base from which other languages could grow. Mumbai trains expanded
my tolerance and culture. A few professors at imperfect universities showed
me the force of good teaching. TIFR showed me serious research. Purdue made
science international and converted independence from an idea into a daily
routine. The lab taught me to remain with a problem that did not care about
my confidence. Teaching sent some knowledge back to where it had started.
Industry connected one experiment to a much larger ambition.

A CV draws this as a smooth upward line. It never felt smooth while I was
inside it. The real education was rebuilding the person doing the learning
every time the environment changed.

\scene

I expected finishing to provide clarity. It mostly gave me a more advanced
version of confusion! The difference was that I no longer needed certainty as
proof that I belonged. A career can change, a company can fail, a field can
turn, and the next phase may again begin in survival mode. After this many
trials, that possibility no longer sounds like evidence of failure.

The lab will stay. Mumbai will continue without waiting. Boston may become a
memory, a home, or only a chapter. Parents will age. Friends will build lives
whose daily details I cannot always share. New people will enter these places,
change them, and leave their own traces. This is not a comfortable thought,
but comfort has never stopped it from being true.

Still, change does not mean that nothing remains. You carry things forward:
the language you once struggled with, a friendship started by complaining on
a flight, the experiment that worked after you nearly abandoned it, the road
you later showed your parents, the habit of waking without an alarm, and the
willingness to begin a conversation.

In the year I turned thirty, I could finally say that my formal education was
ending. It had been long, longer than anyone would aspire for. But one thing
was clearly not ending: the learning. Fortunately or unfortunately, that is
the only truth!

% ══════════════════════════════════════════════════════════════════════
\backmatter
\clearpage
\starchapter{Epilogue: The Latest Version}
% ══════════════════════════════════════════════════════════════════════

\dropcap{I}{n my first weeks in America,} I wrote that there is no way to
document your exact self at one moment. Habits, interests, and emotions keep
changing, so you can access only the latest version and live with it.

This book is my attempt to argue with my own sentence. Very on-brand, I think.

I cannot recover the earlier versions completely, but they left records. One
panics while Mumbai disappears below an aeroplane. One compares every American
habit with India. One nearly fails his courses in a quiet apartment. One hears
his adviser use the word \textit{easier} and decides to stay. One sees a
single atom appear on a camera and cannot sleep properly from excitement.
Another drives his parents from Chicago, walks beside a river in Boston, and
starts calling an internship a job. All of them are me. None is the final
version.

For a long time I assumed life should move toward stability. Now I am not so
sure. No, I am not suggesting changing careers daily or marrying a hundred
people! But destinations become ordinary surprisingly fast. Most of the fun,
frustration, and growth happen while approaching them. Goals are necessary
mainly because they build journeys.

Confusion makes us ask questions. Differences keep conversations alive.
Distance makes certain forms of love visible. Failure tells you what to
measure next. Maybe stability is not the final achievement we imagine. Staying
interested, changing phases, and remaining available to people may be enough.
Really? Not sure. But it is not complete nonsense, right?!

I have stopped trying to subtract what this journey cost from what it gave.
The two came together. The gain was often louder while events happened, and
the loss sometimes became louder later. Both are true.

The experiment continues, and unfortunately the scientist is also the subject.
I observe, participate, lose coherence, recover it, and modify the apparatus
whenever possible. The distance between two homes was not only the thing
separating them. It also built the person travelling between them.

There will be another version after this one. Let us see what he thinks, no?

\textit{One experiment, one conversation, one moment at a time.}

\vfill
\begin{center}
{\color{accent}\rule{1.3in}{0.8pt}}\\[0.3in]
{\small\itshape Everything can be fixed.\\ What cannot be fixed can be lived without.}\\[0.5in]
{\footnotesize\color{soft} West Lafayette, Indiana \(\cdot\) 2026}
\end{center}

% ══════════════════════════════════════════════════════════════════════
\clearpage
\starchapter{About the Author}
% ══════════════════════════════════════════════════════════════════════

\dropcap{S}{aumitra Phatak} grew up in Thane, Maharashtra, and completed his
bachelor's and master's degrees in physics in Mumbai before joining Purdue
University's physics PhD programme in 2021. His doctoral research is in
ultracold atomic physics, the cooling and imaging of single atoms in optical
tweezers, toward the assembly of ultracold polar molecules. He completed a
quantum computing internship in Boston in 2025.

He has been writing about life, science, and the spaces between the two since
2021, in the series \textit{Curious Writings.} He still wakes up without an
alarm.

\vspace{1.4em}
\noindent{\small\color{soft}\url{https://saumitraphatak.github.io/curious-writings}}

\end{document}
